\documentclass{report}
\usepackage{amsmath,amssymb}
\setlength{\parindent}{0mm}
\setlength{\parskip}{1em}
\begin{document}
\begin{center}
\rule{6in}{1pt} \
{\large Matthias Bolten \\
{\bf Advanced smoothing techniques and aggressive coarsening for massively parallel multigrid}}

Universit\"at Kassel \\ Institut f\"ur Mathematik \\ Heinrich-Plett-Stra\ss{}e 40 \\ 34132 Kassel \\ Germany
\\
{\tt bolten@mathematik.uni-kassel.de}\end{center}

Multigrid methods are amongst the most efficient methods for the solution
of partial differential equations in a variety of applications,
especially when the solution of elliptic partial differential equations
is needed. They are used for many simulations in computational science
and engineering, often requiring the use of supercomputers due to the
vast problem size.

While parallel multigrid methods show an optimal scaling behavior on
parallel computers that depends on the number of processors
logarithmically, only, the large number of cores that is available
nowadays on large supercomputers exposes this logarithmic factor. While
this cannot be prevented, the use of aggressive coarsening can reduce
this effect. In order to maintain the efficiency and good algorithmic
scalability of multigrid methods, more powerful smoothing techniques have
to be employed. Different options exist, including polynomial smoothers
or block smoothers that are strongly connected to domain decomposition
methods.

In this talk we will discuss the effect of aggressive coarsening in
geometric multigrid methods and different smoothing techniques that in
combination result in efficient multigrid methods. Analyses and the
necessary techniques will be presented and the advantage of these
techniques on large scale supercomputers and modern computer
architectures will be highlighted.


\end{document}
